\chapter{SDR Drivers}
\label{sec:sdr-drivers}

This chapter describes all of the available drivers for software-defined radios.

\section{Ettus Research}

\begin{tabularx}{16cm}{lllX}
\thickhline
\textbf{Device Family} & \textbf{Driver} & \textbf{Transport} & \textbf{Notes} \\
\thickhline
USRP & uhd & twinlan & \\
\thickhline
\end{tabularx}

\subsection{uhd}

This driver connects via a TCP socket to a socket server
\href{https://github.com/ngscopeclient/scopehal-uhd-bridge}{scopehal-uhd-bridge} which connects to the appropriate
instrument using the UHD API.

This provides network transparency for USB-attached instruments, as well as a license boundary between the BSD-licensed
libscopehal core and the GPL-licensed UHD API.

\section{Analog Devices}

\begin{tabularx}{16cm}{lllX}
\thickhline
\textbf{Device Family} & \textbf{Driver} & \textbf{Transport} & \textbf{Notes} \\
\thickhline
ADALM-PLUTO, other AD9361/AD9363/AD9364 SDRs & iio & iio & Receive, and transmit of tones from the DDS \\
\thickhline
\end{tabularx}

\subsection{iio}

This driver talks to radios based on the AD9361 family (including the AD9363 in the ADALM-PLUTO) using
\href{https://github.com/analogdevicesinc/libiio}{libiio} directly, without a bridge process. It is only available
if ngscopeclient was built with libiio 0.x (0.26 has been tested); newer 1.x releases are not supported yet.

The transport path is a libiio context URI such as \texttt{ip:192.168.2.1} (the default address of a PlutoSDR
connected by USB), \texttt{usb:1.5.5}, or \texttt{local:}. A bare hostname or IP address is treated as an \texttt{ip:}
URI. USB-attached radios are listed in the Endpoint drop-down of the Add Instrument dialog. Radios with two receive
paths show up as two complex channels.

The center frequency sets the receive local oscillator and the span sets the analog RF bandwidth. The sample rate is
configured separately, and should be at least as wide as the span. The limits for each of these come from the radio
(the AD9363 in a stock PlutoSDR cannot tune below 325 MHz, for example, while an AD9361 or a PlutoSDR with unlocked
firmware can), so values you enter are limited to what the radio supports.

Radios with the usual DDS core in the FPGA (the stock firmware of the PlutoSDR and similar radios) can also transmit.
Each transmit path shows up in the stream browser as a channel (TX1, and TX2 on radios with two transmit paths) after
the receive channels, and can generate two tones. Each tone has an on/off switch, a frequency and an amplitude. The
frequency is relative to the transmit local oscillator (LO), so negative frequencies are below the LO, and can be up to
half of the sample rate either side of it. The amplitude is a percentage of the full scale of the DAC, so keep the sum
of the amplitudes of both tones at or below 100\% to avoid clipping. The transmit LO frequency, which is shared by all
transmit paths and is independent of the receive LO, is the TX LO setting in the frequency section of the instrument.
These settings are saved in sessions. The transmitter is not turned off when ngscopeclient disconnects.
The tone amplitudes are not the output power, which also depends on the transmit attenuation of the radio (currently
not adjustable from ngscopeclient).

The receiver gain is in the Input section of the channel properties. The gain mode selects between manual gain and
the radio's automatic gain control (AGC) modes. The gain can only be set by hand in manual mode. A radio that has just
powered up is normally running AGC, and ngscopeclient leaves that alone until you change it. The gain mode and gain
are saved in sessions along with the other settings.

Some things to be aware of:
\begin{itemize}
\item Transmit is limited to tones from the DDS. Arbitrary waveforms cannot be transmitted.
\item Data is captured in blocks of the selected memory depth, so consecutive acquisitions are not gap free.
\item The link to the radio limits the usable sample rate. Over USB 2.0 the higher rates will not keep up.
\end{itemize}

For development and testing without hardware, use the transport path \texttt{mock:} (a simulated single-channel
AD9363) or \texttt{mock:ad9361} (two channels). Transmit settings are accepted by the simulated radio, but nothing
is transmitted. The simulated radio receives a few fixed RF tones in the 433~MHz,
915~MHz and 2.4~GHz bands, which move around as you retune it. In manual gain mode the simulated signal gets stronger
and weaker as you change the gain.

\section{Microphase}

The AntSDR running antsdr\_uhd firmware is supported by the ``uhd" driver for Ettus Research SDRs. There is currently no
support for the IIO firmware.
